\documentclass[11pt]{article}
\usepackage{amsmath,amssymb,amsthm}
\usepackage{hyperref}
\usepackage[margin=1in]{geometry}

\newcommand{\N}{\mathbb{N}}
\newcommand{\Z}{\mathbb{Z}}
\newcommand{\Q}{\mathbb{Q}}
\newcommand{\La}{\Lambda}
\newcommand{\GK}{\operatorname{GK}}
\newcommand{\CPP}{\operatorname{CPP}}
\newcommand{\qpoch}[2]{(#1;#2)_\infty}
\newcommand{\qpochn}[3]{(#1;#2)_{#3}}
\newcommand{\qbin}[2]{\genfrac{[}{]}{0pt}{}{#1}{#2}}
\DeclareMathOperator{\tr}{tr}
\DeclareMathOperator{\wt}{wt}
\DeclareMathOperator{\sgn}{sgn}
\DeclareMathOperator{\Fix}{Fix}
\DeclareMathOperator{\id}{id}

\theoremstyle{plain}
\newtheorem{theorem}{Theorem}
\newtheorem{conjecture}[theorem]{Conjecture}
\newtheorem{proposition}[theorem]{Proposition}
\newtheorem{lemma}[theorem]{Lemma}
\newtheorem{corollary}[theorem]{Corollary}
\theoremstyle{definition}
\newtheorem{definition}[theorem]{Definition}
\newtheorem{example}[theorem]{Example}
\newtheorem{remark}[theorem]{Remark}

\title{The CPP Ground Set \\
  and the $(S, U, \pi)$ Framework}
\author{Robin Langer}
\date{\today}

\begin{document}
\maketitle

\begin{abstract}
We reformulate the ground set for the sign-reversing
involution on~$Q_{n,c}(q)$.  The previous framework
used pairs~$(S, \mathfrak{d})$ where~$\mathfrak{d}$ is
an ascending list of cylindric diagrams (ALCD), which
requires the polynomials~$h_m(q) = (q)_m \cdot g_m(q)$
to have nonneg coefficients.  Since this is not guaranteed
in general, we replace the ALCD by a genuine cylindric
plane partition~$\pi$ and absorb the signs from~$(q)_m$
into an auxiliary subset~$U$.  The resulting ground set
consists of triples~$(S, U, \pi)$ and does not require
any nonnegativity assumption beyond the obvious~$g_m \geq 0$.
We describe the lattice path interpretation: $\pi$~is
always a family of non-intersecting paths on the cylinder,
while~$S$ and~$U$ contribute two independent layers of
alternating signs.
\end{abstract}

\tableofcontents


%% ============================================================
\section{Motivation: why \texorpdfstring{$h_m$}{h\_m} is not enough}
%% ============================================================

Recall the $q$-binomial decomposition
\begin{equation}\label{eq:qbin-decomp}
  Q_{n,c}(q)
  = \sum_{j=0}^{n} (-1)^j \, q^{\binom{j+1}{2}}
    \qbin{n}{j}_{\!q} \; h_{n-j}(q),
\end{equation}
where $h_m(q) = \qpochn{q}{q}{m} \cdot g_m(q)$ and
$g_m(q) = [z^m]\,\GK_c(z,q)$ is the generating function
for cylindric plane partitions of height~$m$.

In the previous framework (CANONICAL\_PEEL.tex), the ground
set was
\[
  \mathcal{G}_{n,c}
  = \bigl\{(S, \mathfrak{d})
    : S \subseteq [n],\;
      \mathfrak{d} \in H_{n-|S|, c}\bigr\},
\]
where~$H_{m,c}$ is a combinatorial set of cardinality
$h_m(q)$ (counted with $q$-weight).  This requires every
coefficient of~$h_m(q)$ to be nonneg.

\begin{remark}
Computationally, $h_m(q)$ is nonneg for all profiles
tested with $d \not\equiv 0 \pmod{3}$ and $d \leq 8$.
However, for $d \equiv 0 \pmod{3}$ it \emph{is} negative.
\end{remark}

\begin{example}[$d = 3$: $h_1$ has negative coefficients]
\label{ex:h1-negative}
For $c = (2, 1, 0)$, $d = 3$, the CPP count at height~$1$
is
\[
  g_1(q) = 2q + 3q^2 + 3q^3 + 3q^4 + 4q^5 + 3q^6
    + 3q^7 + 4q^8 + 3q^9 + 3q^{10} + 4q^{11} + \cdots
\]
The coefficients oscillate with period~$3$:
$\ldots, 4, 3, 3, 4, 3, 3, \ldots$.  Since
$h_1 = (1-q) \cdot g_1$, we get
\[
  h_1(q)
  = 2q + q^2 + q^5 - q^6 + q^8 - q^9 + q^{11} - q^{12}
    + \cdots
\]
The coefficient of~$q^k$ in~$h_1$ is $g_1[k] - g_1[k-1]$,
which is negative whenever~$g_1$ drops (e.g.\
$g_1[6] = 3 < 4 = g_1[5]$, giving $h_1[6] = -1$).

This means the ALCD ground set $\{(S, \mathfrak{d})\}$ has
$-1$ elements at weight~$6$ --- not a genuine set.
\end{example}

The fix is to go back to the \emph{always-nonneg}
quantity~$g_m(q)$, which literally counts CPPs, and to
absorb the alternating signs from~$(q)_m$ into a separate
subset.


%% ============================================================
\section{The new ground set}
%% ============================================================

We start from the identity
\begin{equation}\label{eq:Q-from-gk}
  Q_{n,c}(q)
  = \sum_{j=0}^{n} (-1)^j \, q^{\binom{j+1}{2}}
    \cdot \frac{\qpochn{q}{q}{n}}{\qpochn{q}{q}{j}}
    \cdot g_{n-j}(q).
\end{equation}
This is obtained from~\eqref{eq:qbin-decomp} by writing
$\qbin{n}{j}_q \cdot \qpochn{q}{q}{n-j} = \qpochn{q}{q}{n}/\qpochn{q}{q}{j}$
and noting $h_{n-j} = \qpochn{q}{q}{n-j} \cdot g_{n-j}$.

We decompose each factor combinatorially.

\medskip\noindent
\textbf{Factor 1: the Euler part.}
The term $(-1)^j q^{\binom{j+1}{2}}$ comes from the
initial segment $S = \{1, 2, \ldots, j\}$:
\begin{equation}\label{eq:euler-factor}
  (-1)^j \, q^{1 + 2 + \cdots + j}
  = (-1)^{|S|} \, q^{\sigma(S)},
  \qquad S = [j].
\end{equation}

\medskip\noindent
\textbf{Factor 2: the $(q)_n / (q)_j$ part.}
\begin{equation}\label{eq:qn-qj-factor}
  \frac{\qpochn{q}{q}{n}}{\qpochn{q}{q}{j}}
  = \prod_{k=j+1}^{n} (1 - q^k)
  = \sum_{U \subseteq \{j+1, \ldots, n\}}
    (-1)^{|U|} \, q^{\sigma(U)}.
\end{equation}

\medskip\noindent
\textbf{Factor 3: the CPP part.}
\begin{equation}\label{eq:cpp-factor}
  g_{n-j}(q)
  = \sum_{\pi \in \CPP_{n-j, c}} q^{\wt(\pi)},
\end{equation}
where $\CPP_{m,c}$ denotes the set of rank-$r$ cylindric
plane partitions of profile~$c$ and height~$m$.

\begin{definition}[The ground set $\mathcal{G}_{n,c}$]
\label{def:ground-set}
The \emph{ground set} is
\[
  \mathcal{G}_{n,c}
  = \Bigl\{(S, U, \pi)
    : \begin{array}{l}
        S = \{1, \ldots, j\} \text{ for some } 0 \leq j \leq n, \\
        U \subseteq \{j+1, \ldots, n\}, \\
        \pi \in \CPP_{n - j,\, c}
      \end{array}
    \Bigr\}.
\]
Each element carries:
\begin{itemize}
  \item \textbf{Sign:} $\varepsilon(S, U, \pi) = (-1)^{|S| + |U|}$.
  \item \textbf{Weight:} $w(S, U, \pi) = \sigma(S) + \sigma(U) + \wt(\pi)$.
\end{itemize}
\end{definition}

\begin{proposition}\label{prop:ground-set-identity}
\[
  Q_{n,c}(q)
  = \sum_{(S, U, \pi) \in \mathcal{G}_{n,c}}
    (-1)^{|S| + |U|} \, q^{\sigma(S) + \sigma(U) + \wt(\pi)}.
\]
\end{proposition}

\begin{proof}
Expand~\eqref{eq:Q-from-gk} using
\eqref{eq:euler-factor}--\eqref{eq:cpp-factor}.
\end{proof}


%% ============================================================
\section{The combined set \texorpdfstring{$T = S \cup U$}{T = S ∪ U}}
%% ============================================================

Set $T = S \cup U \subseteq [n]$.  Since $S = [j]$ and
$U \subseteq \{j+1, \ldots, n\}$, the sets $S$ and~$U$
are disjoint and we have
\[
  |T| = |S| + |U| = j + |U|,
  \qquad
  \sigma(T) = \sigma(S) + \sigma(U).
\]
So the sign and weight simplify:
\[
  \varepsilon = (-1)^{|T|},
  \qquad
  w = \sigma(T) + \wt(\pi).
\]

\begin{remark}
The decomposition $T = S \cup U$ is \emph{not} unique: a
subset $T \subseteq [n]$ can arise from different values
of~$j$.  Specifically, $j$ can be any value in
$\{0, 1, \ldots, n\}$ such that $[j] \subseteq T$.  The
maximum such~$j$ is the length of the initial consecutive
run in~$T$:
\[
  j_{\max}(T) = \max\{k \geq 0 : \{1,\ldots,k\} \subseteq T\}.
\]
Different choices of~$j$ give different CPP
heights~$n - j$, so the ground set is genuinely
\emph{three-dimensional}: the triple~$(S, U, \pi)$
carries more information than the pair~$(T, \pi)$.
\end{remark}


%% ============================================================
\section{Lattice path interpretation}
\label{sec:lattice-paths}
%% ============================================================

Each component of the triple~$(S, U, \pi)$ has a clear
role in the lattice path picture.

\subsection*{The CPP $\pi$: non-intersecting paths}

A cylindric plane partition $\pi = (\nu^0, \nu^1,
\ldots, \nu^{r-1})$ of profile~$c$ and height~$m$ encodes
a family of~$r$ non-intersecting lattice paths on the
cylinder with~$m$ levels.  Each level corresponds to an
entry value $1 \leq s \leq m$.  The interlacing conditions
\[
  \nu^i_{p} \geq \nu^{i+1}_{p + c_i}
  \qquad \text{(indices mod $r$)}
\]
ensure that the paths do not cross.

\emph{Key point:} $\pi$ is always a family of
\textbf{non-intersecting} paths.  There are no
``intersecting CPPs'' in this framework.

\subsection*{$S = \{1, \ldots, j\}$: peeled layers}

The initial segment $S = [j]$ records that the
bottom~$j$ layers of a height-$n$ path family have been
\emph{peeled off}.  The CPP $\pi$ has height~$n - j$
(it lives on the \emph{inner} cylinder after removing
the $j$~outermost layers).

Geometrically: a height-$n$ CPP has layers
$1, 2, \ldots, n$.  Layer~$s$ is the horizontal strip
between the $(s{-}1)$-th and $s$-th level.  Peeling
layers $1, \ldots, j$ means these layers are ``detached''
from the CPP and contribute only their level numbers
$1 + 2 + \cdots + j = \sigma(S)$ to the weight.

The sign $(-1)^j$ reflects that each peeled layer
introduces one ``virtual crossing'' between the
detached paths and the remaining family.

\subsection*{$U \subseteq \{j+1, \ldots, n\}$: base removal}

In the E--H factorisation
$\GK_c = F_c / \qpoch{zq}{q}$, the factor
$1/\qpochn{q}{q}{m} = \sum_\lambda q^{|\lambda|}$ (sum
over partitions with parts~$\leq m$) generates a
\emph{base partition}: an ordinary partition that slides
underneath all the lattice paths without affecting their
relative positions.  Every CPP of height~$m$ contains
this trivial base.

The factor $\qpochn{q}{q}{m}$ in
$h_m = \qpochn{q}{q}{m} \cdot g_m$ \emph{removes} the
base partition via inclusion--exclusion.  Expanding
$\prod_{k=j+1}^n (1 - q^k)$: each factor $(1 - q^k)$
offers a binary choice --- leave level~$k$ alone
(coefficient~$1$) or subtract its base strip
(coefficient~$-q^k$).  The subset $U$ records which
levels are being subtracted.

\subsection*{Two types of defect}

The triple $(S, U, \pi)$ carries two independent layers
of sign:

\medskip
\begin{center}
\begin{tabular}{l|l|l}
  & \textbf{$S$ (crossings)} &
    \textbf{$U$ (base removal)} \\ \hline
  \textbf{What} &
    peeled layers &
    subtracted base strips \\
  \textbf{Range} &
    $S = \{1, \ldots, j\}$ &
    $U \subseteq \{j{+}1, \ldots, n\}$ \\
  \textbf{NILP meaning} &
    virtual crossings between &
    inclusion--exclusion for \\
  & detached and remaining paths &
    removing trivial padding \\
  \textbf{Sign} &
    $(-1)^{|S|}$ &
    $(-1)^{|U|}$ \\
  \textbf{Weight} &
    $\sigma(S) = 1 + \cdots + j$ &
    $\sigma(U)$
\end{tabular}
\end{center}

\medskip\noindent
The total sign is $(-1)^{|S|+|U|}$ and the total weight
is $\sigma(S) + \sigma(U) + \wt(\pi)$.  Fixed points of
the involution must have $|S| + |U|$ even; the ideal fixed
points are $(\emptyset, \emptyset, \pi)$ with sign~$+1$.


%% ============================================================
\section{The involution}
\label{sec:involution}
%% ============================================================

To prove positivity of~$Q_{n,c}(q)$, we need a
sign-reversing, weight-preserving involution
$\iota \colon \mathcal{G}_{n,c} \to \mathcal{G}_{n,c}$
whose fixed points all have positive sign.

The sign is $(-1)^{|S| + |U|}$, so $\iota$ must change
the parity of~$|S| + |U|$.  There are three types of
elementary moves.

\subsection{Type 1: flip a \texorpdfstring{$U$}{U}-element}

Choose some $k \in \{j+1, \ldots, n\}$ and toggle its
membership in~$U$:
\[
  (S, U, \pi) \;\longmapsto\; (S, U \triangle \{k\}, \pi').
\]
Here $\pi$ does not change, so weight preservation
requires $\sigma(U \triangle \{k\}) = \sigma(U) \pm k$
and $\wt(\pi')$ adjusts accordingly.

If $\pi' = \pi$ (the CPP is unchanged), then the weight
changes by $\pm k$ from the $U$-part alone.  This means a
pure $U$-flip can only be weight-preserving if we also
modify $\pi$.

\subsection{Type 2: peel/glue}

Extend or shrink the initial segment:
\[
  (S, U, \pi) \;\longmapsto\; (S', U', \pi'),
  \qquad |S'| = |S| \pm 1.
\]
Peeling ($j \to j+1$): detach layer $j+1$ from $\pi$,
producing a CPP~$\pi'$ of height~$n - j - 1$.
The weight of the detached layer (horizontal strip at
level $j+1$) must be absorbed.

Gluing ($j \to j-1$): attach a new outermost layer
to~$\pi$, producing $\pi'$ of height $n - j + 1$.

\subsection{Type 3: transfer at the boundary}

Move element $j+1$ between $S$ and~$U$:
\begin{itemize}
  \item If $j+1 \in U$: remove $j+1$ from~$U$ and
    extend $S$ to $[j+1]$.  This increases~$|S|$ by~1
    and decreases~$|U|$ by~1, so $|S|+|U|$ is unchanged
    --- \textbf{no sign change}.  Not useful for the
    involution directly, but useful as a preparatory step.
  \item If $j+1 \notin U$: extend $S$ to $[j+1]$ and
    the parity of $|S|+|U|$ increases by~1.  This \emph{is}
    a sign-reversing move.
\end{itemize}


%% ============================================================
\section{Comparison with the ALCD framework}
\label{sec:comparison}
%% ============================================================

\begin{center}
\begin{tabular}{lll}
  \hline
  & \textbf{ALCD framework} & \textbf{CPP framework} \\
  \hline
  Ground set & $(S, \mathfrak{d})$
    & $(S, U, \pi)$ \\[2pt]
  $S$ & $S \subseteq [n]$ (arbitrary)
    & $S = [j]$ (initial segment) \\[2pt]
  Combinatorial object
    & $\mathfrak{d} \in H_{m,c}$ (ALCD)
    & $\pi \in \CPP_{m,c}$ (CPP) \\[2pt]
  Requires & $h_m(q) \geq 0$
    & $g_m(q) \geq 0$ (automatic) \\[2pt]
  Sign & $(-1)^{|S|}$
    & $(-1)^{|S| + |U|}$ \\[2pt]
  Weight & $\sigma(S) + \wt(\mathfrak{d})$
    & $\sigma(S) + \sigma(U) + \wt(\pi)$ \\[2pt]
  NILP meaning & $\mathfrak{d}$ encodes
    & $\pi$ is always \\
  & (non)-intersecting paths
    & non-intersecting \\[2pt]
  Involution & peel/glue at one level
    & three types of moves \\
  \hline
\end{tabular}
\end{center}

The key advantage of the CPP framework is robustness:
$g_m(q)$ is a nonneg polynomial by definition (it counts
CPPs), so the ground set is well-defined without any
additional assumptions.

The cost is an extra layer of structure: the subset~$U$
adds a second source of alternating signs.  But this also
provides more flexibility for the involution, since there
are more elementary moves available.


%% ============================================================
\section{Structure of the involution}
\label{sec:structure}
%% ============================================================

The involution $\iota$ on $\mathcal{G}_{n,c}$ should
proceed as follows.

\begin{enumerate}
  \item \textbf{If $U \neq \emptyset$:}
    try to cancel using a $U$-flip.
    Let $k = \min(U)$ or $k = \max(U)$ (the canonical
    choice).  If there exists $\pi'$ with
    $\wt(\pi') = \wt(\pi) + k$ (or $\wt(\pi) - k$) and
    the same height, pair
    $(S, U, \pi) \leftrightarrow (S, U \setminus \{k\}, \pi')$.

  \item \textbf{If $U = \emptyset$ and $j \geq 1$:}
    try to cancel using a peel/glue.
    The outermost layer of $\pi$ at level $j$ can be
    detached (if the strip size matches) or a new layer
    can be attached.

  \item \textbf{If $U = \emptyset$ and $j = 0$:}
    the element $(S, U, \pi) = (\emptyset, \emptyset, \pi)$
    has sign $+1$.  This is a candidate \textbf{fixed point}.
    It contributes to~$Q_{n,c}(q)$ if and only if no
    valid pairing exists.
\end{enumerate}

The fixed points of~$\iota$ are triples
$(\emptyset, \emptyset, \pi)$ that survive all
cancellation attempts.  These are CPPs of height~$n$
that contribute $q^{\wt(\pi)}$ to~$Q_{n,c}(q)$.

\begin{remark}
The fixed points are a \emph{subset} of the CPPs of
height~$n$.  Not every height-$n$ CPP survives: some
are paired with elements having $U \neq \emptyset$ or
$S \neq \emptyset$.  The surviving CPPs are exactly
those that generate~$Q_{n,c}(q)$.
\end{remark}


%% ============================================================
\section{Three strategies for the involution}
\label{sec:strategies}
%% ============================================================

The two signed sets $S$ and $U$ can be handled in several
ways.

\subsection{Strategy 1: sequential sieve}

First cancel all $U$-signs (the $(q)_m$ inclusion--exclusion),
then cancel all $S$-signs (the crossing cancellation).
This is equivalent to:
\begin{enumerate}
  \item Prove $h_m \in \N[q]$ (the $U$-sieve produces
    a nonneg result), then
  \item apply the Franklin involution from
    CANONICAL\_PEEL.tex to the residual.
\end{enumerate}
This is the original approach.  It works when
$h_m \geq 0$, but fails for $d \equiv 0 \pmod{3}$ because
the $U$-sieve does not close up cleanly.

\subsection{Strategy 2: interleaved sieve}

At each step, choose whether to toggle an element of $S$
or an element of $U$, based on a priority rule.  The
Franklin comparison generalises: compare the ``cheapest''
$U$-flip (smallest $k \in U$ or smallest $k \notin U$) with
the cheapest peel/glue, and do whichever is smaller.

This interleaves the two types of cancellation, potentially
resolving cases where the sequential approach gets stuck.

\subsection{Strategy 3: combined set}

Define $T = S \cup U \subseteq [n]$ and run a single
Franklin sieve on~$T$.  The sign is $(-1)^{|T|}$ and the
weight is $\sigma(T) + \wt(\pi)$.  At each level~$k$,
toggle $k \in T$ (adding or removing~$k$), adjusting the
CPP height if $k$ is at the $S$--$U$ boundary.

The difficulty: toggling $k \in T$ may change whether $k$
belongs to $S$ or $U$, which changes the CPP height, which
changes the set of available CPPs.  The coupling between
$|S|$ and the CPP height is the main obstruction.


%% ============================================================
\section{The priority rule}
\label{sec:priority}
%% ============================================================

The involution must choose which move to apply.
We propose the following priority:

\begin{enumerate}
  \item \textbf{$U$-cancellation first.}
    If $U \neq \emptyset$, work with~$U$ before
    touching~$S$ or~$\pi$.  The $U$-flip is the simplest
    move (it does not change the CPP height), so it
    should be preferred.

  \item \textbf{Smallest element.}
    When cancelling via~$U$, always target
    $k = \min(U)$ (analogous to the Franklin
    involution targeting the smallest part).

  \item \textbf{Peel/glue as fallback.}
    Only modify~$S$ (and hence the CPP height) when
    no $U$-cancellation is available.
\end{enumerate}

This priority ensures:
\begin{itemize}
  \item Elements with $U \neq \emptyset$ are handled
    first, leaving only $(S, \emptyset, \pi)$ elements
    for the peel/glue step.
  \item The $U$-cancellation step is entirely within the
    set of CPPs of the \emph{same height}, making it
    easier to verify involutivity.
  \item The peel/glue step operates on a smaller residual
    set, reducing the complexity of the canonical peel
    problem.
\end{itemize}


%% ============================================================
\section{Open questions}
%% ============================================================

\begin{enumerate}
  \item \textbf{$U$-cancellation mechanism.}
    When $k \in U$, how exactly is the partner
    CPP~$\pi'$ with $\wt(\pi') = \wt(\pi) \pm k$
    chosen?  This requires a canonical bijection
    between weight-$w$ and weight-$(w+k)$ CPPs of the
    same height, which exists when $g_m$ has
    nondecreasing coefficients (up to a point) but
    needs to be made explicit.

  \item \textbf{Interaction between $U$ and peel/glue.}
    After the $U$-cancellation step eliminates all
    elements with $U \neq \emptyset$, the residual
    elements are $(S, \emptyset, \pi)$.  These have
    sign $(-1)^{|S|}$ and weight
    $\sigma(S) + \wt(\pi)$.  This is exactly the ALCD
    framework restricted to initial-segment subsets~$S$.
    The peel/glue involution from
    CANONICAL\_PEEL.tex applies to this residual.

  \item \textbf{Does the $U$-step close up?}
    For the $U$-cancellation to define an involution
    on its own, we need: if $(S, U, \pi)$ maps to
    $(S, U', \pi')$, then applying the same rule to
    $(S, U', \pi')$ returns to $(S, U, \pi)$.  This
    requires the bijection $\pi \leftrightarrow \pi'$
    to be its own inverse.

  \item \textbf{Weight constraint.}
    The $U$-flip at level~$k$ changes the $U$-weight
    by~$\pm k$.  For weight preservation, $\wt(\pi)$
    must change by~$\mp k$.  This means we need to
    find, for each CPP~$\pi$ of height~$m$, a
    partner~$\pi'$ of the same height with
    $\wt(\pi') = \wt(\pi) \pm k$.  This is a matching
    problem on CPPs --- nontrivial but concrete.

  \item \textbf{Fixed points of the $U$-step.}
    Elements $(S, U, \pi)$ with $U \neq \emptyset$ that
    cannot be $U$-cancelled (no valid $\pi'$ exists)
    must be handled by a different mechanism.  Do these
    exist?  If so, they need to be routed to the
    peel/glue step.
\end{enumerate}


%% ============================================================
\section{Connection to previous notes}
%% ============================================================

\begin{itemize}
  \item \textbf{CANONICAL\_PEEL.tex:}
    The peel/glue framework for the ALCD ground set.
    Sections 3--4 (Franklin comparison, weight condition)
    apply directly to the residual after $U$-cancellation.

  \item \textbf{Q\_NILP.tex:}
    The enlarged path space and crossing framework.
    Section~4.3 (intersecting paths as relaxed interlacing)
    describes what happens when $S \neq \emptyset$.

  \item \textbf{Q\_LOCAL.tex:}
    The local rule and conflict detection ($A \cap B$).
    Relevant for the peel/glue step on the residual.
\end{itemize}


\end{document}
